\section{Comparison, implementation and further questions}
\label{sec:comparison}
\subsection{Stationary positive realization and finite causal width}
Benvenuti--Farina \cite[Theorems 2 and 3]{BF} formulate existence of
positive realizations by invariant polyhedral cones and shift-invariant
positive sequences. Their Example 4 uses an irrational planar rotation
to prevent a finite stationary realization. Monr\`as--Winter
\cite[Theorem 6 and Lemma 7]{MW} give the corresponding stochastic
cone criterion and a rotation example. These are direct antecedents,
not merely neighboring subjects.

Theorem~\ref{thm:headline} makes a different finite assertion. The
matrices are allowed to be unrelated at different epochs, and a new
machine may be selected at every horizon. Each entire finite array is
positively realizable, while each of its cuts separately needs exactly
three states. The lower bound is the accumulation of enlargement along
$P_t\subset P_{t+1}$ and $RP_t\subset P_{t+1}$. An invariant-cone
existence theorem or a peripheral-spectrum argument for one common
matrix does not, by itself, give this quantitative inequality for
changing rows. Conversely, our theorem does not re-prove or extend the
full stationary positive-realization existence classification.

The reachable-section construction is deliberately compatible with
classical quotient realizations. In \cite[Section 4, Theorem 9]{MW},
the accessible subspace and its unobservable kernel determine the
minimal linear quotient; arbitrary positive realizations can contain
more directions. That is why Lemma~\ref{lem:lift} uses a section of a
simplex before projecting. It does not identify minimum hidden dimension
with a minimum number of vertices of an unlifted observable polygon.
The related distinction between restricted and unrestricted positive
factorizations appears explicitly in \cite[Definition 1, Theorem 1]{GG}.
The exponential allowance $2^{K_t}$ is conservative, but makes the
lower bound valid for all competing hidden representations.

The recent work of Lumbreras--Ma--Thompson--Gu \cite{LMTG} studies
irreducible errors and decision failures of finite classical world
models, with exact quantum realizations. Its Section IV.2 and ensuing
proofs use repetition of one update, Perron--Frobenius limits and an
irrational phase. Our task has no quantum simulator or decision-loss
claim: it compares exact \emph{finite-horizon} cut-dependent classical
realizations. The common rotation ancestry and different row-sharing
quantifiers are both relevant. We do not infer a novelty boundary from
the absence of a citation to this manuscript in that work.

\subsection{Residual restrictions, direct sums and sampling}
The finite-generation and canonical residual constructions of
Esposito--Lemay--Denis--Dupont \cite{ELDD} and Denis--Esposito
\cite[Proposition 16]{DE} concern state languages constrained by the
residual representation. The lower bound here places no such constraint
on a state's full continuation. Indeed all separate minima stay three
in the rotation example, so a large list of actual residuals alone would
not prove its causal width lower bound. Classical hidden-Markov and
filtered-experiment comparison remain part of the surrounding theory
\cite{Vidyasagar,Norberg}; our wordwise finite profile statement is not
claimed to be a new general comparison theorem for filtered spaces.

Nonnegative direct-sum additivity itself is familiar, including the
tensor version in \cite[Lemma 4.1]{QCL}. Theorem~\ref{thm:tagged-sum}
uses the same positivity principle to separate \emph{reachable state
sets at every time}; it identifies a sum of entire feasible profile
sets. Taking only sums of separate ranks would lose the tradeoff in
\eqref{eq:growing-frontier}. Its tag and full common suffix alphabet are
explicit and cannot be omitted.

The additive encoder construction in the preserved article \cite{GTF31}
uses weighted replacement. The sampling step is weighted reservoir
sampling, as described, for example, in \cite[Appendix A.2]{PBRT}:
replace the current sample with probability equal to the new weight
divided by total weight. When total previous weight is zero, the dummy
label is chosen from the existing alphabet. Neither that sampling
algorithm nor the classical simplex inputs are novelty claims of the
present paper. The earlier online consequences and rational sampler
remain available without being needed for the rotation lower bound.

\subsection{What is and is not charged}
The finite-horizon bounds concern exact atomic stochastic transducers,
not a bit-operation model. Both alphabet cardinality and all persistent
hidden choices are charged. The row table and the prescribed external
epoch are not. The rational family permits separately charged fair-bit
sampling; the algebraic or computable-real polygon rows are not given a
finite-denominator implementation claim. Sampling and verification costs
therefore do not disappear into the phrase ``exact realization.''

The new lower bound holds even with the external clock and arbitrary
cut-dependent rows. The polygon upper construction has common command
matrices across its training epochs, but its prescribed decoder depends
on $N$. This does not define one finite machine that works at every
horizon with an unchanged decoder. No automatic transfer from the
clocked peak to autonomous length recognition is made.

\subsection{Conclusions}
The proved separation is between \emph{separate positive minima} and
\emph{simultaneous positive width}, with fixed alphabets, uniformly
positive response probabilities, and an arbitrarily long horizon.
The projected reachable-section budget applies to a class of
volume-preserving mean dynamics. The golden-angle family has an
explicit multi-cut inequality and a constructive square-root upper
bound; the rational family has effective excluded-horizon certificates.
The tagged operation separately supplies exact growing Pareto frontiers.

The optimal growth of $W_N$, a sharp measure of lift complexity,
quantitative stability uniform in the horizon, and general optimization
of informative acquisition or validation actions remain distinct
questions. The logarithmic lower bound is not matched by the
square-root upper bound. Complete arbitrary-profile minimization is not
claimed. These distinctions do not affect the unbounded separation or
the exact tagged frontier.

The historical geometric, large-deviation and nonlinear-operator chains
are preserved as separate mathematics. No finite-array argument here is
claimed to discharge their analytic gates. The focused article contains
the proofs required for its new claims; the cumulative volumes are
archival records, not an additional burden of certification for its
referee.
